\documentclass[11pt]{article}

\usepackage[utf8]{inputenc}
\usepackage[T1]{fontenc}
\usepackage{amsmath,amssymb,amsthm}
\usepackage{mathtools}
\usepackage{hyperref}
\usepackage{enumitem}
\usepackage[margin=1in]{geometry}

\newtheorem{theorem}{Theorem}[section]
\newtheorem{lemma}[theorem]{Lemma}
\newtheorem{proposition}[theorem]{Proposition}
\newtheorem{corollary}[theorem]{Corollary}
\theoremstyle{definition}
\newtheorem{definition}[theorem]{Definition}
\newtheorem{example}[theorem]{Example}

\hypersetup{
    colorlinks=true,
    linkcolor=blue,
    citecolor=blue,
    urlcolor=blue
}

\title{Exact Optimal Values for Quantum Local Packing via the Resultant Method}

\author{Euzebio Santos\\
\textit{Independent Researcher, Brazil}\\
\texttt{ebiossanto@github.com}}

\date{September 2026}

\begin{document}

\maketitle

\begin{abstract}
We prove that for any graph $G = (V,E)$ and any dimension $d \in \mathbb{N}$,
the optimal value $c_{\mathrm{opt}}(G,d)$ of the quantum local packing problem
is a root of an explicit polynomial $P_G(c) \in \mathbb{Z}[c]$, computable via
the resultant method applied to the rank constraints of the Gram matrix.
We give a complete proof, explicit computations for small graphs, and
tight degree bounds.
\end{abstract}

\section{Introduction}

\subsection{Problem Statement}

Given a graph $G = (V,E)$ with $|V| = n$ and dimension $d \in \mathbb{N}$,
the \emph{quantum local packing problem} asks for:
\begin{equation}\label{eq:problem}
c_{\mathrm{opt}}(G,d) = \min \left\{ c \in \mathbb{R} \,:\, \exists\, n_1,\ldots,n_n \in \mathbb{R}^d \text{ unit},\; \langle n_i, n_j \rangle \le c\ \forall\, (i,j) \in E \right\}
\end{equation}

This problem arises in quantum contextuality~\cite{klyachko2008, pirónio2010},
quantum key distribution bounds, and entanglement theory.

\subsection{Prior Work}

Known results include:
\begin{itemize}
    \item Complete graphs $K_n$: $c_{\mathrm{opt}} = -1/(n-1)$ (regular simplex)
    \item Lov\'{a}sz theta: $\vartheta(G)$ provides a semidefinite upper bound
    \item For $K_5$ minus edge in $\mathbb{R}^3$: $c_{\mathrm{opt}}$ is the root
    of a cubic~\cite{soares2026}
\end{itemize}

\subsection{Our Contribution}

We prove a general theorem: for arbitrary $G$ and $d$, the optimal value
$c_{\mathrm{opt}}(G,d)$ is algebraic, and its minimal polynomial is computable
via the resultant method.

\section{Preliminaries}

\begin{definition}[Gram Matrix]
For $n$ unit vectors $n_1, \ldots, n_n \in \mathbb{R}^d$, the \emph{Gram matrix}
is $G_{ij} = \langle n_i, n_j \rangle$.
\end{definition}

The Gram matrix satisfies:
\begin{enumerate}[label=(\roman*)]
    \item $G = G^T$ (symmetric)
    \item $G_{ii} = 1$ for all $i$ (unit vectors)
    \item $G \succeq 0$ (positive semidefinite)
    \item $\operatorname{rank}(G) \le d$
\end{enumerate}

\begin{theorem}[Rank and Minors]\label{thm:rank}
$\operatorname{rank}(G) \le d$ if and only if all $(d+1) \times (d+1)$ minors
of $G$ are zero.
\end{theorem}

This is a standard result in linear algebra. The key consequence is that
the rank constraint is equivalent to $\binom{n}{d+1}$ polynomial equations
of degree $d+1$ in the off-diagonal entries of $G$.

\section{Main Result}

\begin{theorem}[Existence of Characteristic Polynomial]\label{thm:main}
Let $G = (V,E)$ be a graph with $|V| = n$ and let $d \in \mathbb{N}$.
Then:
\begin{enumerate}
    \item There exists a polynomial $P_G(c) \in \mathbb{Z}[c]$ such that
    $c_{\mathrm{opt}}(G,d)$ is a root of $P_G$.
    \item $P_G$ is computable via the resultant method: eliminate the free
    Gram entries from the minor equations $M_S(G) = 0$.
    \item The degree satisfies:
    \begin{equation}\label{eq:bound}
    \deg(P_G) \le (d+1)^{\binom{n}{d+1}}
    \end{equation}
\end{enumerate}
\end{theorem}

\begin{proof}
\textbf{Step 1: Polynomial system.}

Let $A \subseteq E$ be the set of active edges at the optimum
(where $G_{ij} = c$). Let $F = \binom{n}{2} - |A|$ be the number of
free off-diagonal entries.

The system is:
\begin{align}
M_S(G) &= 0 \quad \forall\, S \subseteq V,\ |S| = d+1 \label{eq:minors}\\
G_{ii} &= 1 \quad \forall\, i \label{eq:diag}\\
G_{ij} &= c \quad \forall\, (i,j) \in A \label{eq:active}
\end{align}

This is a system of $\binom{n}{d+1}$ polynomial equations of degree $d+1$
in $F$ free variables and the parameter $c$.

\textbf{Step 2: Elimination via resultant.}

We eliminate the free variables $y_1, \ldots, y_F$ one at a time:
\begin{enumerate}
    \item Choose two equations $f_i, f_j$ that involve $y_F$
    \item Compute $R_1 = \operatorname{Res}_{y_F}(f_i, f_j)$ --- a polynomial
    in $c, y_1, \ldots, y_{F-1}$
    \item Replace the system with $\{R_1\} \cup \{f_k : y_F \notin f_k\}$
    \item Repeat for $y_{F-1}, \ldots, y_1$
\end{enumerate}

After $F$ steps, we obtain a polynomial $P_G(c)$ in $c$ alone.

\textbf{Step 3: Degree bound.}

Each resultant $\operatorname{Res}_y(f,g)$ has degree
$\deg_y(f) \cdot \deg_y(g)$ in the remaining variables. Since each $M_S$
has degree $d+1$ in each variable:
\begin{equation}
\deg(P_G) \le (d+1)^{\binom{n}{d+1}}
\end{equation}

This is B\'{e}zout's bound. For sparse graphs, free variables reduce the
effective system size, giving tighter bounds.
\end{proof}

\begin{proposition}[Feasibility Filter]\label{prop:filter}
A root $c_0$ of $P_G(c)$ is feasible (i.e., $c_0 \ge c_{\mathrm{opt}}$)
if and only if:
\begin{enumerate}
    \item There exists a Gram matrix $G$ with $G_{ii} = 1$,
    $G_{ij} \le c_0$ for all $(i,j) \in E$
    \item $G \succeq 0$
    \item $\operatorname{rank}(G) \le d$
\end{enumerate}

The optimal value is $c_{\mathrm{opt}} = \min\{c_0 : c_0 \text{ is a feasible root}\}$.
\end{proposition}

\section{Explicit Computations}

\subsection{$K_4$ in $\mathbb{R}^2$}

\begin{theorem}
$c_{\mathrm{opt}}(K_4, 2) = -1/3$.
\end{theorem}

\begin{proof}
The Gram matrix of the regular simplex in $\mathbb{R}^2$ is
$G = (1+1/3)I - (1/3)J$, with all off-diagonal entries $= -1/3$.

The characteristic polynomial is:
\begin{equation}
P_{K_4}(c) = (1+3c)(1-c)^3
\end{equation}

Roots: $c = -1/3$ (triple) and $c = 1$ (triple).

At $c = -1/3$: eigenvalues are $\{0, 4/3, 4/3, 4/3\}$, rank $= 2$.
Feasible.
\end{proof}

\subsection{$K_5$ minus edge in $\mathbb{R}^3$}

\begin{theorem}\label{thm:k5m}
For $q = 49/50$, $c_{\mathrm{opt}}(K_5 - \{04\}, 3) = c_* \approx -0.3046$,
where $c_*$ is the physical root of:
\begin{equation}
P(c,q) = 9c^3 - (q+2)c^2 - (2q+3)c - q = 0
\end{equation}
\end{theorem}

\begin{proof}
The proof combines:
\begin{enumerate}
    \item \textbf{Algebraic identity:} $\det(G) = -F_1 \cdot F_2$
    \item \textbf{Resultant:} Eliminating the free variable $d = u \cdot v$
    gives $4(q-c)(c-1)^2 \cdot P(c,q) = 0$.
    \item \textbf{Rank analysis:} root$_0$ has Gram eigenvalue $< 0$ (infeasible);
    root$_2 > 0$ (not minimum).
    \item \textbf{Primal certificate:} explicit configuration with rank 3.
\end{enumerate}
\end{proof}

\subsection{$K_5$ in $\mathbb{R}^2$}

\begin{theorem}
$c_{\mathrm{opt}}(K_5, 2) = \cos(2\pi/5) = \frac{\sqrt{5}-1}{4} \approx 0.309$.
\end{theorem}

\begin{proof}
Regular pentagon with vertices at angles $0, 2\pi/5, 4\pi/5, 6\pi/5, 8\pi/5$.
All pairwise dot products $\le \cos(2\pi/5)$. Rank 2, PSD.
\end{proof}

\subsection{Cycle $C_4$ in $\mathbb{R}^2$}

\begin{theorem}
$c_{\mathrm{opt}}(C_4, 2) = -1/2$.
\end{theorem}

\begin{proof}
The characteristic polynomial from the resultant method is:
\begin{equation}
P_{C_4}(c) = 8c^4 - 8c^2 - c + 1 = (c-1)(2c+1)(2c^2 - 2c - 1)
\end{equation}

Roots: $c = 1, -1/2, (1 \pm \sqrt{3})/2$.

At $c = -1/2$: the symmetric Gram matrix with $G_{03}$ free is PSD with
rank 2. The root $c = (1-\sqrt{3})/2 \approx -0.366$ is infeasible
(edge $(0,3)$ violates the constraint).
\end{proof}

\subsection{Path $P_4$ in $\mathbb{R}^2$}

\begin{theorem}
$c_{\mathrm{opt}}(P_4, 2) = -1$.
\end{theorem}

\begin{proof}
Alternating vectors: $n_0 = (+1, 0)$, $n_1 = (-1, 0)$, $n_2 = (+1, 0)$,
$n_3 = (-1, 0)$. All edge dot products $= -1$. Rank 1, PSD.
\end{proof}

\section{Degree Bounds}

\begin{theorem}[B\'{e}zout Bound]\label{thm:bezout}
For the complete graph $K_n$ in dimension $d$:
\begin{equation}
\deg(P_{K_n}) \le (d+1)^{\binom{n}{d+1}}
\end{equation}
\end{theorem}

\begin{table}[h]
\centering
\begin{tabular}{ccccl}
\hline
$n$ & $d$ & B\'{e}zout bound & Actual degree & Case \\
\hline
3 & 1 & $2^3 = 8$ & 1 & Trivial \\
4 & 2 & $3^4 = 81$ & 3 & Regular simplex \\
5 & 2 & $3^5 = 243$ & --- & $\cos(2\pi/5)$ \\
5 & 3 & $4^5 = 1024$ & 3 & Cubic \\
6 & 3 & $4^{20} \approx 10^{12}$ & ? & Open \\
\hline
\end{tabular}
\caption{Degree bounds for complete graphs.}
\end{table}

\begin{theorem}[Sparsity Reduction]
For a graph $G$ with $|E|$ edges, the effective number of free variables is
$F = \binom{n}{2} - |E|$. The degree bound improves to:
\begin{equation}
\deg(P_G) \le (d+1)^{\binom{n}{d+1} - F}
\end{equation}
\end{theorem}

\section{Algorithm}

\begin{quote}
\textbf{Algorithm} \textsc{ResultantPacking}

\textbf{Input:} Graph $G = (V,E)$, dimension $d$

\textbf{Output:} Polynomial $P_G(c)$, optimal value $c_{\mathrm{opt}}$

\begin{enumerate}
    \item Construct symbolic Gram matrix $\mathbf{G}$ with entries $G_{ij}$
    \item Set $G_{ii} = 1$ for all $i$
    \item For each edge $(i,j) \in E$: set $G_{ij} = c$ (parameter)
    \item Compute all $\binom{n}{d+1}$ minors $M_S = \det(\mathbf{G}[S,S])$
    for $|S| = d+1$
    \item \textbf{Eliminate free variables:}
    \begin{itemize}
        \item While free variables remain:
        \item Pick a free variable $y$
        \item Pick $f, g$ involving $y$
        \item Compute $R = \operatorname{Res}_y(f, g)$
        \item Replace $f, g$ with $R$
    \end{itemize}
    \item $P_G(c) = \prod_{f \in \mathcal{F}} f$ (or GCD)
    \item Find all real roots of $P_G(c)$
    \item For each root $c_0$, check feasibility (PSD + rank $\le d$)
    \item Return $c_{\mathrm{opt}} = \min\{c_0 : c_0 \text{ feasible}\}$
\end{enumerate}
\end{quote}

\section{Comparison with SDP}

\begin{table}[h]
\centering
\begin{tabular}{lccc}
\hline
Method & $K_5-\{04\}/\mathbb{R}^3$ & Exact? & Complexity \\
\hline
SDP (Lov\'{a}sz) & $c = -0.4714$ & No (upper bound) & Polynomial \\
Resultant & $c_* = -0.3046$ & \textbf{Yes} & Exponential \\
NLP & $c_* = -0.3046$ & Numerical only & Slow \\
\hline
\end{tabular}
\caption{Comparison of methods for $K_5$ minus edge.}
\end{table}

\section{Open Problems}

\begin{enumerate}
    \item \textbf{Complexity:} Is computing $P_G(c)$ polynomial-time for fixed $d$?

    \item \textbf{Tight bounds:} What is the exact degree of $P_G$ for complete graphs?

    \item \textbf{Graph families:} Do trees, planar graphs, or regular graphs have
    special structure that simplifies $P_G$?

    \item \textbf{Anti-commutativity:} Can this framework handle fermionic
    (anti-commuting) systems?

    \item \textbf{Higher dimensions:} What happens for $d \ge 4$?
\end{enumerate}

\begin{thebibliography}{9}

\bibitem{klyachko2008}
A.~Klyachko, M.~A.~Can, M.~Binicio{\u{g}}lu, and A.~S.~Shumovsky,
``Simple test for hidden variables in spin-1 systems,''
\textit{Phys. Rev. Lett.} \textbf{101}, 020403 (2008).

\bibitem{pirónio2010}
S.~Pironio, M.~Navascu\'{e}s, and A.~Ac\'{i}n,
``Convergent relaxations of polynomial optimization problems with
non-commuting variables,''
\textit{SIAM J. Optim.} \textbf{20}, 2157--2180 (2010).

\bibitem{lovasz1979}
L.~Lov\'{a}sz,
``On the Shannon capacity of a graph,''
\textit{IEEE Trans. Inf. Theory} \textbf{25}, 1--7 (1979).

\bibitem{soares2026}
E.~Soares,
``Quantum Local Packing: K5 Minus Edge in R$^3$,''
GitHub repository (2026).

\end{thebibliography}

\end{document}
